\documentclass[10pt,a4paper]{amsart}
\usepackage{amsmath,amssymb,amsthm,graphicx}
% Retained minimal preamble for the Stacks Project chapter sites-cohomology.tex (Cohomology on Sites).
% Verbatim excerpt of the chapter preamble: Xy-pic setup (native preamble lines 14-18) and the native
% theorem declarations with the mathematical macros (native preamble lines 154-217). Copied mechanically;
% nothing added, removed or redefined. The wrapper class and the amsmath/amssymb/amsthm/graphicx base
% packages are supplied by the assembler, so the class/externaldocument part of the native preamble is omitted.
% For commutative diagrams we use Xy-pic
\usepackage[all]{xy}

% We use 2cell for 2-commutative diagrams.
\xyoption{2cell}
% Theorem environments.
%
\theoremstyle{plain}
\newtheorem{theorem}[subsection]{Theorem}
\newtheorem{proposition}[subsection]{Proposition}
\newtheorem{lemma}[subsection]{Lemma}

\theoremstyle{definition}
\newtheorem{definition}[subsection]{Definition}
\newtheorem{example}[subsection]{Example}
\newtheorem{exercise}[subsection]{Exercise}
\newtheorem{situation}[subsection]{Situation}

\theoremstyle{remark}
\newtheorem{remark}[subsection]{Remark}
\newtheorem{remarks}[subsection]{Remarks}

\numberwithin{equation}{subsection}

% Macros
%
\def\lim{\mathop{\mathrm{lim}}\nolimits}
\def\colim{\mathop{\mathrm{colim}}\nolimits}
\def\Spec{\mathop{\mathrm{Spec}}}
\def\Hom{\mathop{\mathrm{Hom}}\nolimits}
\def\Ext{\mathop{\mathrm{Ext}}\nolimits}
\def\SheafHom{\mathop{\mathcal{H}\!\mathit{om}}\nolimits}
\def\SheafExt{\mathop{\mathcal{E}\!\mathit{xt}}\nolimits}
\def\Sch{\mathit{Sch}}
\def\Mor{\mathop{\mathrm{Mor}}\nolimits}
\def\Ob{\mathop{\mathrm{Ob}}\nolimits}
\def\Sh{\mathop{\mathit{Sh}}\nolimits}
\def\NL{\mathop{N\!L}\nolimits}
\def\CH{\mathop{\mathrm{CH}}\nolimits}
\def\proetale{{pro\text{-}\acute{e}tale}}
\def\etale{{\acute{e}tale}}
\def\QCoh{\mathit{QCoh}}
\def\Ker{\mathop{\mathrm{Ker}}}
\def\Im{\mathop{\mathrm{Im}}}
\def\Coker{\mathop{\mathrm{Coker}}}
\def\Coim{\mathop{\mathrm{Coim}}}

% Boxtimes
%
\DeclareMathSymbol{\boxtimes}{\mathbin}{AMSa}{"02}

%
% Macros for moduli stacks/spaces
%
\def\QCohstack{\mathcal{QC}\!\mathit{oh}}
\def\Cohstack{\mathcal{C}\!\mathit{oh}}
\def\Spacesstack{\mathcal{S}\!\mathit{paces}}
\def\Quotfunctor{\mathrm{Quot}}
\def\Hilbfunctor{\mathrm{Hilb}}
\def\Curvesstack{\mathcal{C}\!\mathit{urves}}
\def\Polarizedstack{\mathcal{P}\!\mathit{olarized}}
\def\Complexesstack{\mathcal{C}\!\mathit{omplexes}}
% \Pic is the operator that assigns to X its picard group, usage \Pic(X)
% \Picardstack_{X/B} denotes the Picard stack of X over B
% \Picardfunctor_{X/B} denotes the Picard functor of X over B
\def\Pic{\mathop{\mathrm{Pic}}\nolimits}
\def\Picardstack{\mathcal{P}\!\mathit{ic}}
\def\Picardfunctor{\mathrm{Pic}}
\def\Deformationcategory{\mathcal{D}\!\mathit{ef}}
\title{A variant of the derived projection formula for a morphism of ringed topoi}
\author{The Stacks Project Authors}
\date{stacks-project:sites-cohomology.tex:local-snapshot-20261011; GNU FDL 1.2 or later}
\begin{document}
\maketitle
\noindent\textit{Editorial note.} All delivered mathematical text is the author's own, copied from the pinned Stacks Project chapter file sites-cohomology.tex; no statement or proof step was rewritten, reordered or removed.\par\medskip
\noindent\textit{Editorial note.} The proof is the author's own argument. It displays the module-level identity and then states, in the author's own words, that the lemma follows from this and the definitions; that closing sentence is retained verbatim inside the counted proof span and is not replaced or expanded by this card.\par\medskip
\noindent\textit{Editorial note.} The counted spans contain no label reference and no bibliography citation, so the card carries no navigation stub and no bibliography excerpt.\par\medskip
\noindent\textit{Editorial note.} Cross-chapter restatement: the same label occurs in chapter cohomology for a ringed space. That statement is the special case of this one, so per the lane-s4 rule only this more general occurrence is counted here.\par\medskip
\noindent\textit{Editorial note.} The preamble is the chapter's own retained minimal preamble (sources/stacks-sites-cohomology/retained-preamble.tex), a verbatim excerpt of the native preamble: the Xy-pic setup and the native theorem declarations with the mathematical macros.\par\medskip
\noindent\textit{Allowed standard prerequisites (editorial).}\begin{itemize}
\item Morphism of ringed topoi and the functors f\textasciicircum{}\{-1\}, f\textasciicircum{}* and Lf\textasciicircum{}*. Same chapter, the inverse image f\textasciicircum{}\{-1\} of sheaves and the derived pullback Lf\textasciicircum{}* on modules over a ringed topos. Supplies the functors appearing on both sides of the asserted isomorphism.
\item Derived tensor product of complexes of modules over a ringed topos. Same chapter, the derived tensor product tensor\textasciicircum{}L and its computation from K-flat representatives. Supplies the objects compared and the passage from the module-level identity to complexes.
\item The identification f\textasciicircum{}*G = O tensor over f\textasciicircum{}\{-1\}O' with f\textasciicircum{}\{-1\}G. Same chapter, the description of the pullback of a module as a tensor product over the inverse image of the structure sheaf. Supplies the displayed identity whose proof is the whole content of the argument.
\end{itemize}\medskip
\begin{lemma}
\label{lemma-variant-derived-pullback}
Let $f : (\Sh(\mathcal{C}), \mathcal{O}) \to (\Sh(\mathcal{C}'), \mathcal{O}')$
be a morphism of ringed topoi. There is a canonical bifunctorial
isomorphism
$$
\mathcal{F}^\bullet
\otimes_\mathcal{O}^{\mathbf{L}}
Lf^*\mathcal{G}^\bullet
=
\mathcal{F}^\bullet 
\otimes_{f^{-1}\mathcal{O}'}^{\mathbf{L}}
f^{-1}\mathcal{G}^\bullet 
$$
for $\mathcal{F}^\bullet$ in $D(\mathcal{O})$ and
$\mathcal{G}^\bullet$ in $D(\mathcal{O}')$.
\end{lemma}

\begin{proof}
Let $\mathcal{F}$ be an $\mathcal{O}$-module and let $\mathcal{G}$
be an $\mathcal{O}'$-module. Then
$\mathcal{F} \otimes_{\mathcal{O}} f^*\mathcal{G} =
\mathcal{F} \otimes_{f^{-1}\mathcal{O}'} f^{-1}\mathcal{G}$
because
$f^*\mathcal{G} =
\mathcal{O} \otimes_{f^{-1}\mathcal{O}'} f^{-1}\mathcal{G}$.
The lemma follows from this and the definitions.
\end{proof}

\end{document}
